% Related Work Section
% This file contains the background and related work

\section{Related Work}
\label{sec:related_work}

\paragraph{Positional Encoding in Transformers.}
Transformers \cite{vaswaniAttentionAllYou2023} lack inherent positional awareness due to the permutation-equivariant nature of self-attention. To address this, the original Transformer used sinusoidal absolute position embeddings 
(APE), which add fixed position-dependent signals to input tokens. Learnable APE \cite{devlinBERTPretrainingDeep2019} replaces fixed sinusoids 
with trainable vectors, offering greater flexibility. However, APE methods encode absolute positions, which limits their ability to generalize to 
sequence lengths unseen during training. Relative position encodings address this limitation by encoding the distance between tokens rather than 
their absolute positions. Shaw et al.~\cite{shawSelfAttentionRelativePosition2018} introduced learnable relative position biases added to 
attention logits. T5 \cite{raffelExploringLimitsTransfer2023} further simplified this approach with a small number of learnable bias buckets. 
ALiBi \cite{pressTrainShortTest2022} proposed a non-learned linear bias that scales with distance, enabling length extrapolation in language 
models. When extending Transformers from 1D sequences to 2D visual data, vision Transformer (ViT) \cite{dosovitskiyImageWorth16x162021} applies Transformers to images by dividing them into patches, using learnable APE. DeiT \cite{touvronTrainingDataefficientImage2021} enabled efficient ViT training on ImageNet-1K, while Swin Transformer \cite{liuSwinTransformerHierarchical2021} introduced relative position bias within local windows. These developments highlight the importance of positional encoding design in vision models for handling varying resolutions and capturing spatial relationships.

\paragraph{Rotary Position Embedding.}
Rotary Position Embedding (RoPE) \cite{suRoFormerEnhancedTransformer2023} introduces a fundamentally different approach to position encoding by applying rotation operations directly to query and key vectors. Unlike additive position encodings, RoPE multiplies the query and key representations by rotation matrices whose angles are determined by the position index and a set of base frequencies. This design ensures that the dot product between query and key depends only on their relative positions, inherently encoding relative positional information. RoPE has become the de facto position encoding for modern large language models including LLaMA \cite{touvronLLaMAOpenEfficient2023} and Qwen \cite{yangQwen3TechnicalReport2025}, demonstrating excellent length extrapolation capabilities. Various extensions have been proposed to improve RoPE's ability to handle long sequences, including position interpolation \cite{chenExtendingContextWindow2023} which rescales positional indices at inference time and YaRN \cite{pengYaRNEfficientContext2023} which refines this strategy through frequency-aware scaling to improve stability.

\paragraph{2D RoPE for Vision.}
Extending RoPE to 2D vision data requires careful design choices. The straightforward approach, which we term axial 2D RoPE, splits the embedding dimension into two halves and applies 1D RoPE independently along the x and y axes. The 2D RoPE method has been adopted in many vision tasks, such as SAM 2 \cite{raviSAM2Segment2024} for image segmentation and Qwen-Image \cite{wuQwenImageTechnicalReport2025} for image generation. M-RoPE (Multimodal RoPE) \cite{wangQwen2VLEnhancingVisionLanguage2024} decomposes image indices into height, width, and temporal components to align visual patches with 1D text sequences. However, these axial approaches encode only horizontal and vertical relationships, without explicit awareness of diagonal or other directional spatial patterns. Heo et al.~\cite{heoRotaryPositionEmbedding2024} proposed RoPE-Mixed, treats the frequencies for both spatial axes as learnable parameters, allowing the network to adaptively "mix" axial signals to capture diagonal information. In multimodal settings, several works generalize RoPE to jointly handle tokens from different modalities. For instance, Circle-RoPE \cite{wangCircleRoPEConelikeDecoupled2025} project image token indices onto a ring that is orthogonal to the linear axis of text token indices to eliminate spurious cross-modal biases.
